 %%%%%%%%%%%%%%%%%%%%%%%%%%%%%%%%%%%%%%%%%
% "ModernCV" CV and Cover Letter
% LaTeX Template
% Version 1.11 (19/6/14)
%
% This template has been downloaded from:
% http://www.LaTeXTemplates.com
%
% Original author:
% Xavier Danaux (xdanaux@gmail.com)
%
% License:
% CC BY-NC-SA 3.0 (http://creativecommons.org/licenses/by-nc-sa/3.0/)
%
% Important note:
% This template requires the moderncv.cls and .sty files to be in the same 
% directory as this .tex file. These files provide the resume style and themes 
% used for structuring the document.
%
%%%%%%%%%%%%%%%%%%%%%%%%%%%%%%%%%%%%%%%%%


%----------------------------------------------------------------------------------------
%	PACKAGES AND OTHER DOCUMENT CONFIGURATIONS
%----------------------------------------------------------------------------------------

\documentclass[12pt,a4paper,sans]{moderncv} % Font sizes: 10, 11, or 12; paper sizes: a4paper, letterpaper, a5paper, legalpaper, executivepaper or landscape; font families: sans or roman

\usepackage{fontawesome5}
\usepackage{xcolor}


\moderncvcolor{red} % CV color - options include: 'blue' (default), 'orange', 'green', 'red', 'purple', 'grey' and 'black'
\moderncvstyle{classic} % CV theme - options include: 'casual' (default), 'classic', 'oldstyle' and 'banking'
\definecolor{venetianred}{HTML}{A42A04} 
\colorlet{color1}{venetianred} % main accent color \colorlet{color2}{mycolor!70} % secondary (slightly lighter)


%\colorlet{color2}{mycolor!70} % secondary (slightly lighter)

%\definecolor{color1}{HTML}{434C98}

\usepackage{geometry}
 \geometry{
 a4paper,
 total={210mm,297mm},
 left=10mm,
 right=10mm,
 top=20mm,
 bottom=20mm,
 }


%----------------------------------------------------------------------------------------
%	NAME AND CONTACT INFORMATION SECTION
%----------------------------------------------------------------------------------------

\firstname{Marco} % Your first name
\familyname{Bellò} % Your last name


\mobile{+393497617003}
 % Phone number
\email{marco.bello.vi@gmail.com} % Email ID



%-------------------------------------------
%%%%%%  RESUME STARTS HERE  %%%%%%%%%%%%%%%%%%%%%%%%%%%%

\begin{document}

%----------HEADING----------
\begin{center}
    \textbf{\Huge Marco Bellò} \\ \vspace{5pt}
    \small \faMobile* +39 3497617003 \hspace{1pt} $|$
    \hspace{1pt} \faEnvelope \hspace{2pt} \href{mailto:marco.bello.vi@gmail.com}{marco.bello.vi@gmail.com} \hspace{1pt} $|$ 
    \hspace{1pt} \faGithub \hspace{2pt} \href{https://github.com/mhetacc}{github.com/mhetacc} \hspace{1pt} $|$
    \hspace{1pt} \faTelegram \hspace{2pt} \href{https://t.me/mhetac}{t.me/mhetac}
    \\ \vspace{-3pt}
\end{center}
\vspace{10pt} % adjust spacing as needed

{\color{venetianred} \hrule height 0.8pt}

\vspace{10pt}
%----------------------------------------------------------------------------------------
%	EDUCATION SECTION
%----------------------------------------------------------------------------------------
\section{Interests and Relevance}

I am a Master's student in Computer Science at the University of Padua. I will graduate in the second half of September 2026.

Starting from a security and networking track, I kept finding myself leaning towards data analysis and machine learning-related projects. Since I always had a passion for linguistics, discourse, and communication, especially when connected to political and social contexts, I spent the last year specializing in the field of Artificial Intelligence, with special regard to linguistic and social applications. 

\vspace{2mm}

I have extensive experience with the Python programming language, since I used it for both ``common'' applications, such as data visualization and machine learning, and for less common ones: I simulated a concurrent online-gaming scenario leveraging remote procedure calls and multithreading.

I also used the R programming language to perform classification and regression tasks, and to study the presence of echo-chambers in social networks and the degree of their segregation.

\vspace{2mm}

\textbf{My background and research interests directly align with the project's objectives and directions:} the Master's Thesis I am working on (more details in the next section) requires me to do both data collection and data analysis using natural language processing tools and pipelines, and has a focus on linguistic itself. Moreover, since the corpus will be in Italian, I will already have experience working with a language that is foreign to most large language models.

I am compiling a corpus of political speeches, which are often ambiguous and open to different interpretations. This makes them a relevant case study since they represent high levels of variation and non-standardness compared to both common and technical speech, on which models are often trained (for example, the whole of Wikipedia). 
%----------------------------------------------------------------------------------------
%	PROJECT SECTION
%----------------------------------------------------------------------------------------

\section{Master Thesis}
 
I am gathering a corpus of speeches from Italian Prime Ministers between 1945 and 2025. I will then analyze textual complexity and classify hate speech and populism, and cross-compare the results with democracy level indicators taken from the V-Dem dataset.

\vspace{2mm}

\textbf{The first step} will be building a corpus of speeches. To limit the scope, I decided to take only speeches from Italian Prime Ministers. The dataset will have a fixed, structured taxonomy inspired by ParlaMint.

It is likely that for older politicians there are fewer speeches available, i.e., the dataset becomes more sparse as we go back in time. For this reason, data augmentation techniques may be useful.

For video or audio speeches, speech to text processing will be performed. I will likely use OpenAI's Whisper library, a Transformer sequence-to-sequence model trained on various speech processing tasks. 

\vspace{2mm}

\textbf{The second step} is using the corpus to perform linguistic analyses. The corpus will be pre-processed, with procedures such as tokenization, stemming and lemmatization. I will then compute part-of-speech tagging, as well as word frequency and text complexity measures, following with hate speech and populism classification tasks. The corpus will have to be translated into vectorial representation.

All of the above can be done with popular Python libraries such as Stanza, spaCy, FastText, LLM2Vec, and textacy, and with transformers-only libraries from Hugging Face. 

\vspace{2mm}

\textbf{The last step} is cross-referencing the results with democratic indicators such as freedom of expression, freedom of association, transparency and fairness of law and other such metrics. All the democracy parameters can be found in the V-Dem dataset.





%----------------------------------------------------------------------------------------
%	ACHIEVEMENTS SECTION
%----------------------------------------------------------------------------------------


%----------------------------------------------------------------------------------------
%	PROFILE% SECTION
%---------------------------------------------------------------------------------------%-%

%\section{Profiles}

%\cvitem{GitHub}{\url{https://github.com/sakshib}}
%\cvitem{clubVA}{\url{https://www.clubva.com/public/employee/avatar/index/Sakshi}\newline{}}

\end{document}
